% ===================================================================
%  LENTELĖ Nr. 3 — Vaikystė ir jaunystė.
%  Traduction 2026 d'apres texto/io, controlee sur texto/fr, texto/ru
%  et texto/cs. Consigne : voir texto/lt/10-lentele-01.tex.
%
%  LE GENITIF ANTEPOSE MORD ENFIN, et il mord exactement la ou le
%  tableau 1 avait annonce qu'il mordrait. Deux inversions, et les
%  deux sont la meme construction :
%
%    * « la flecho (7) dil klosheyo (5) » devient « varpinės (5)
%      smailę (7) », le clocher passe devant sa fleche ;
%    * « la abeli (22) dil abeluyo (21) » devient « avilio (21)
%      bitutės (22) », la ruche passe devant ses abeilles.
%
%  ON NE LES A PAS EVITEES, ET C'EST DELIBERE. On aurait pu tourner
%  chaque fois par une relative — « la fleche, qui couronne le
%  clocher » — et rendre le compte a zero ; mais ce serait ecrire un
%  lituanien de traducteur pour flatter un controle. Le releve compte
%  les inversions, il ne demande pas qu'il n'y en ait pas. Le
%  finnois en avait treize, le basque vingt-neuf, et personne n'a
%  tordu leur syntaxe pour les effacer.
%
%  UN MOT LA OU L'IDO EN A DEUX, ET LE NOMBRE FAIT LA DIFFERENCE.
%  L'alinea 2 met cote a cote « la vetero » et « l'aero » ; le
%  lituanien n'a qu'« oras » pour les deux — mais il le decline au
%  PLURIEL pour le temps qu'il fait et au SINGULIER pour l'air qu'on
%  respire. « Orai gražūs, oras tyras. » Aucune des treize colonnes
%  precedentes n'avait separe deux notions par le seul nombre ; le
%  tcheque avait « počasí » et « vzduch », le galicien « tempo » et
%  « aire ». Ici la distinction ne tient pas au lexique mais a la
%  flexion, et le gras la porte telle quelle.
%
%  LE FAC-SIMILE EMPLOIE DEUX FOIS LE MEME NUMERO — (19) pour le
%  presbytere et pour les abricotiers, (45) pour le taureau et pour
%  le mat de cocagne. C'est de l'imprimeur, non du transcripteur, et
%  les treize colonnes precedentes l'ont toutes garde tel quel.
% ===================================================================

%%K t03-apar-1 apar
\VUsaut{3.0mm}
\VUtitre{15.0pt}{60}{LENTELĖ Nr. 3}
\VUsaut{1.6mm}
\VUtitre{11.4pt}{0}{Vaikystė ir jaunystė.}
\VUsaut{3.4mm}

%%K t03-apar-2 apar
(3-ioji ir 4-oji lentelės sudėtos taip, kad keturiose \VUgras{šeimos scenose}
pateiktų esmines sąvokas apie \VUgras{orus}, \VUgras{amžių}, \VUgras{šeimą},
\VUgras{aprangą} ir \VUgras{padėtį}.)

\VUsaut{4.0mm}
%%K t03-tit-1 sub
\VUcentre{12.0pt}{0}{\textit{Pirmoji scena.}}
\VUsaut{3.0mm}

%%K t03-tit-2 sub
\VUpk{10.2pt}{40}{Krikštynos kaime.}
\VUsaut{2.4mm}

%%K t03-tit-3 sub
\VUpk{10.2pt}{40}{Pavasaris.}
\VUsaut{3.0mm}

%%K t03-c1-01-1 p
1. --- Esame \VUgras{pavasarį,} \VUgras{kovo}, \VUgras{balandžio} ar \VUgras{gegužės}
mėnesį, \VUgras{savaitės} dienomis: \VUgras{pirmadienį}, \VUgras{antradienį} ar
\VUgras{trečiadienį}. \VUgras{Dešimta valanda} ryto.

%%K t03-c1-02-1 p
2. --- \VUgras{Orai} gražūs, \VUgras{oras} tyras. Saulė šviečia. Keli
\VUgras{debesėliai} \textsuperscript{(2)} kyla mėlyname \VUgras{danguje}
\textsuperscript{(1)}.

%%K t03-c1-03-1 p
3. --- \VUgras{Medžiai} vos turi \VUgras{lapų.} Lieknos \VUgras{tuopos}
\textsuperscript{(3)} ir aukštas \VUgras{platanas} \textsuperscript{(17)} kol kas turi
tik pumpurus.

%%K t03-c1-04-1 p
4. --- \VUgras{Kregždės} \textsuperscript{(4)} grįžta ir džiaugsmingai čirendamos
skraido aplink \VUgras{smailę} \textsuperscript{(7)}, kylančią nuo
\VUgras{varpinės} \textsuperscript{(5)}, kuri vainikuoja kaimo \VUgras{bažnyčią}
\textsuperscript{(6)}.

%%K t03-tit-4 sub
\VUsaut{3.4mm}
\VUpk{10.2pt}{40}{Bažnyčia.}
\VUsaut{3.0mm}

%%K t03-c1-05-1 p
5. --- Labai paprasta ta bažnyčia su savo dideliu \VUgras{portalu}
\textsuperscript{(13)} ir stambiu \VUgras{laikrodžiu} \textsuperscript{(8)}. Mažoji
\VUgras{rodyklė} \textsuperscript{(9)} ties skaitmeniu X: ji rodo \VUgras{valandas.}
\VUgras{Didžioji} \textsuperscript{(10)} ties XII (vidudienis): ji rodo
\VUgras{minutes.}

%%K t03-c1-06-1 p
6. --- Varpinėje linksmai skamba \VUgras{varpai} \textsuperscript{(11)}. Stogo
viršūnėje stovi mažas akmeninis \VUgras{kryžius} \textsuperscript{(12)}.

%%K t03-c1-07-1 p
7. --- Bažnyčia turi ne tik didįjį \VUgras{portalą} \textsuperscript{(13)}, ji turi
ir šonines dureles. Pro tas dureles ponas \VUgras{klebonas} \textsuperscript{(15)},
apsivilkęs savo \VUgras{sutana,} išeina iš \VUgras{zakristijos}
\textsuperscript{(14)} ir eina į \VUgras{kleboniją} \textsuperscript{(19)}.

%%K t03-c1-08-1 p
8. --- Šalia jo štai \VUgras{pienininkė} \textsuperscript{(24)} su savo
\VUgras{asilu} \textsuperscript{(23)}. Ji grįžta iš miesto, kur nunešė gero pieno
vaikams.

%%K t03-tit-5 sub
\VUsaut{4.0mm}
\VUpk{10.2pt}{40}{Vaismedžių sodas.}
\VUsaut{3.4mm}

%%K t03-c1-09-1 p
9. --- Ponas Aliboronas (asilas) visiškai patenkintas tuo rytiniu bėgimu. Jis
gardžiai traukia orą, kvepiantį \VUgras{žiedų} \textsuperscript{(20)} kvapu, kurie
dengia gretimo \VUgras{vaismedžių sodo} \VUgras{vaismedžius} \textsuperscript{(18)}.
Visai rausvi ir balti tie \VUgras{persikmedžiai} \textsuperscript{(18)} ir
\VUgras{abrikosmedžiai} \textsuperscript{(19)}, ir jie žada gerą derlių.

%%K t03-c1-10-1 p
10. --- Tuo tarpu mažosios \VUgras{bitutės} \textsuperscript{(22)} iš \VUgras{avilio}
\textsuperscript{(21)} zvimbdamos eina ten grobio rinkti savo saldaus
\VUgras{medaus.}

%%K t03-c1-11-1 p
11. --- Bet kas gi darosi? Štai kaimo vaikai atbėga ir išbaido išsigandusias
\VUgras{žąsis} \textsuperscript{(25)}. Tai palyda išeina iš bažnyčios po notaro sūnaus
\VUgras{krikšto.}

%%K t03-c1-12-1 p
12. --- Mažasis Jokūbas savo \VUgras{lopšyje} \textsuperscript{(26)} pabunda nuo
džiaugsmingo varpų skambėjimo, o jo sesuo Magdalena, kuri irgi nori pamatyti
krikštynų palydą, prašo motiną ateiti iššukuoti jos plaukų ir aprengti ją namų
slenkstyje.

%%K t03-c1-13-1 p
13. --- Motina sutinka. Ji užsideda savo \VUgras{skarelę} \textsuperscript{(28)}, savo
\VUgras{liemenę} \textsuperscript{(29)} ir savo \VUgras{prijuostę} \textsuperscript{(30)}
ir ateina atsisėsti ant mažos kėdutės prieš duris. Magdalena teturi savo
\VUgras{apatinį sijoną} \textsuperscript{(31)} ir savo \VUgras{korsetą}
\textsuperscript{(32)}.

%%K t03-c1-14-1 p
14. --- Kokia ji laiminga! Ji pamiršta savo mėgstamas gėles, savo
\VUgras{gvazdikus} \textsuperscript{(34)}, ant kurių tupia \VUgras{drugeliai}
\textsuperscript{(35)}, savo \VUgras{tulpes} \textsuperscript{(36)}, savo
\VUgras{pakalnutes} \textsuperscript{(37)} \VUgras{vazonuose} \textsuperscript{(38)} ant
\VUgras{sienos} \textsuperscript{(33)}, ir savo \VUgras{rožes} \textsuperscript{(39)},
kurios sudaro tokią gražią gyvatvorę. Ji apžiūrinėja turtingą palydos ponių aprangą.

%%K t03-c1-15-1 p
15. --- Kaimo vaikai ne itin žiūri į aprangą. Jie skuba ir stumdo vieni kitus, kad
gautų \VUgras{saldainių} \textsuperscript{(55)} ar \VUgras{monetų}
\textsuperscript{(52)}. Vienas iš jų verkia \textsuperscript{(40)}.

%%K t03-c1-16-1 p
16. --- Kitas užmynė ant kojos \VUgras{šuniui} \textsuperscript{(42)}, kuris
cypdamas bėga ir išbaido \VUgras{vištą} \textsuperscript{(43)} su jos
\VUgras{viščiukais} \textsuperscript{(44)}.

%%K t03-tit-6 sub
\VUsaut{5.0mm}
\VUpk{10.2pt}{40}{Krikštynų eisena.}
\VUsaut{4.0mm}

%%K t03-c1-17-1 p
17. --- Priešakyje eina \VUgras{auklė} \textsuperscript{(45)}. Ji turi gražią
\VUgras{kyką} \textsuperscript{(46)} su ilgais kaspinais. Ji neša \VUgras{naujagimį}
\textsuperscript{(47)}, visą apvilktą \VUgras{nėriniais} \textsuperscript{(48)}.

%%K t03-c1-18-1 p
18. --- Paskui auklę eina \VUgras{krikštatėvis} \textsuperscript{(49)} ir
\VUgras{krikštamotė} \textsuperscript{(53)}. Jie dalija saldainius ir monetas.
Krikštatėvis mūvi \VUgras{pirštines} \textsuperscript{(51)} ir dėvi
\VUgras{cilindrą} \textsuperscript{(50)}. Krikštamotė rankoje laiko gražią
\VUgras{puokštę} \textsuperscript{(54)}.

%%K t03-c1-19-1 p
19. --- Paskui juos matome vaiko \VUgras{dėdę} \textsuperscript{(56)} ir
\VUgras{tetą} \textsuperscript{(58)}. Teta turi elegantišką \VUgras{skrybėlę}
\textsuperscript{(59)}, dėdė --- juodą \VUgras{redingotą} \textsuperscript{(57)} ir
baltą liemenę. Štai paskui \VUgras{pusbrolis} \textsuperscript{(60)} ir
\VUgras{pusseserė} \textsuperscript{(61)}. Ši ranka laiko naujagimio \VUgras{sesutę}
\textsuperscript{(62)}, mažąją Bertą.

%%K t03-c1-20-1 p
20. --- Kitas Bertos \VUgras{brolis} \textsuperscript{(63)} eina paskutinis. Jis
duoda ranką \VUgras{giminaitei} \textsuperscript{(64)}. Šeimos \VUgras{draugai}
\textsuperscript{(65)} eina paskui.

%%K t03-tit-7 sub
\VUsaut{4.6mm}
\VUpk{10.2pt}{40}{Žiūrovai.}
\VUsaut{3.4mm}

%%K t03-c1-21-1 p
21. --- Prie palydos štai \VUgras{elgeta} \textsuperscript{(66)}, atsirėmęs į savo
\VUgras{ramentą} \textsuperscript{(67)}. Jis prašo išmaldos (elgetauja).

%%K t03-c1-22-1 p
22. --- Kitoje pusėje yra mažas labai \VUgras{mandagus} \textsuperscript{(68)}
berniukas. Jis sveikinasi ir linki „\VUgras{labos dienos}“ pažįstamiems žmonėms.

%%K t03-c1-23-1 p
23. --- Kaimiečiai išėjo iš savo namų pažiūrėti krikštynų eisenos. Matome
\VUgras{kaimietį} \textsuperscript{(69)} su savo \VUgras{medvilnine kepuraite,} o
šalia \VUgras{kaimietę} \textsuperscript{(70)}. Paskui \VUgras{senutę}
\textsuperscript{(71)}, atsirėmusią į savo \VUgras{lazdą} \textsuperscript{(72)}.
Galiausiai \VUgras{malūnininką} \textsuperscript{(73)} su savo didele \VUgras{veltine
skrybėle} \textsuperscript{(74)} ir jauną kaimietę, apsiavusią \VUgras{klumpėmis}
\textsuperscript{(75)}, kuri moko savo mažylį vaikščioti.

%%K t03-c1-24-1 p
24. --- Ta scena labai linksma.

\VUsaut{4.0mm}
%%K t03-tit-8 sub
\VUcentre{12.0pt}{0}{\textit{Antroji scena.}}
\VUsaut{3.0mm}

%%K t03-tit-9 sub
\VUpk{10.2pt}{40}{Viešoji šventė.}
\VUsaut{2.4mm}

%%K t03-tit-10 sub
\VUpk{10.2pt}{40}{Vandens žaidimai ir regatos.}
\VUsaut{3.0mm}

%%K t03-c2-01-1 p
1. --- Atėjo \VUgras{vasara.} Karšta \VUgras{birželio} (liepos ar
\VUgras{rugpjūčio}) mėnesio saulė skatina jaunuolius maudytis ir irstytis.

%%K t03-c2-02-1 p
2. --- Dešiniajame upės krante irkluotojai rengiasi, kad dalyvautų vandens žaidimuose
ir regatose.

%%K t03-c2-03-1 p
3. --- Vienas iš jų nusiauna savo \VUgras{batus} \textsuperscript{(1)}; jis
atriša (atsagsto) juos. Du jo draugai jau pasirengę. Nuo šiol jie tebeturi tik savo
\VUgras{trikotažinius marškinėlius} \textsuperscript{(2)} ir \VUgras{maudymosi
kelnaites} \textsuperscript{(3)}. Jie šnekučiuojasi laukdami savo draugų. Jie kalba
apie atlaidus ir šventę, kuri vyksta kitame krante.

%%K t03-c2-04-1 p
4. --- Ant žemės šalia jų guli jų draugų \VUgras{drabužiai}: \VUgras{pora}
\VUgras{batelių} \textsuperscript{(4)}, \VUgras{batai} \textsuperscript{(5)},
\VUgras{kojinės} \textsuperscript{(6)}, \VUgras{veltinės} \textsuperscript{(7)} ir
\VUgras{šiaudinės} \textsuperscript{(8)} skrybėlės ir \VUgras{kaklaraištis}
\textsuperscript{(9)}.

%%K t03-c2-05-1 p
5. --- Vienas jaunuolis rūpestingai sulankstė savo \VUgras{švarką}
\textsuperscript{(10)}. Jis deda jį ant rankšluosčio, į kurį jo kaimynas netrukus
suvynios abu \VUgras{drabužius}.

%%K t03-c2-06-1 p
6. --- Šeštasis vaikinas vėluoja. Jis dar tebeturi ant galvos savo \VUgras{kepurę}
\textsuperscript{(11)}. Jis nenusivilko nei savo \VUgras{marškinių}
\textsuperscript{(12)}, nei savo \VUgras{apykaklės} \textsuperscript{(13)}, nei savo
\VUgras{rankogalių} \textsuperscript{(14)}. Rankoje jis laiko savo \VUgras{liemenę}
\textsuperscript{(17)} ir tebeturi savo \VUgras{kelnes} \textsuperscript{(16)} ir
\VUgras{petnešas} \textsuperscript{(15)}. Tikrai jis nebus pasirengęs artimiausioms
lenktynėms.

%%K t03-c2-07-1 p
7. --- Šalia jaunuolių takelis, kurio pakraščiuose daug \VUgras{dagių}
\textsuperscript{(18)}, veda prie upės kranto, kur tėvas Jokūbas tarp \VUgras{nendrių}
\textsuperscript{(19)} rengia \VUgras{fejerverką} \textsuperscript{(20)}, kurį
paleis vakare.

%%K t03-tit-11 sub
\VUsaut{4.4mm}
\VUpk{10.2pt}{40}{Lenktynės.}
\VUsaut{3.4mm}

%%K t03-c2-08-1 p
8. --- Upe kelios ponios, prisidengusios savo skėčiais nuo saulės, plaukioja
\VUgras{valtimi} \textsuperscript{(21)}. Jos seka lenktynes, kurios labai įdomios.
Kiekvienoje \VUgras{jolėje} \textsuperscript{(22)} keturi \VUgras{irkluotojai}
\textsuperscript{(24)} iš visų jėgų iriasi, o \VUgras{vairininkas}
\textsuperscript{(26)} laiko \VUgras{vairą} \textsuperscript{(25)} ir kreipia trapią
valtį. \VUgras{Irklai} \textsuperscript{(23)} ritmingai kerta vandenį, ir lengvos
valtys plaukia svaiginamu greičiu.

%%K t03-c2-09-1 p
9. --- Keli jaunuoliai varžosi prie tilto atramos dėl \VUgras{įstrižosios karties}
\textsuperscript{(28)} prizo. Vienas iš jų atsargiai eina lanksčia ir slidžia
kartimi, kad pasiektų mažą \VUgras{vėliavėlę} \textsuperscript{(29)}, pritvirtintą
gale. Jo nelaimingas \VUgras{varžovas} \textsuperscript{(30)} įkrito į vandenį. Jis
plaukia į krantą.

%%K t03-c2-10-1 p
10. --- Vyresni jaunuoliai \VUgras{rengia valčių turnyrą} \textsuperscript{(32)}
prie \VUgras{krantinės} \textsuperscript{(33)}. Visus tuos žaidimus stebi gausi
\VUgras{publika} \textsuperscript{(31)}, atsirėmusi į \VUgras{tilto}
\textsuperscript{(27)} \VUgras{turėklus} ir susigrūdusi \VUgras{krante}. Vis dėlto
keli žiūrovai atsisuka sekti \VUgras{prizinės karties} \textsuperscript{(45)} žaidimo
arba pasigrožėti \VUgras{mero palyda}, kuri ateina premijų \VUgras{dalyti}.

%%K t03-c2-11-1 p
11. --- Pirmiausia štai keli \VUgras{berniūkščiai} \textsuperscript{(35)}, einantys
pirma \VUgras{muzikantų} \textsuperscript{(36)}; paskui ateina \VUgras{meras}
\textsuperscript{(37)}, padėjėjai ir miestelio \VUgras{savivaldybės taryba.}
Galiausiai žygiuoja \VUgras{ugniagesiai} \textsuperscript{(38)}.

%%K t03-tit-12 sub
\VUsaut{5.0mm}
\VUpk{10.2pt}{40}{Mugė.}
\VUsaut{3.6mm}

%%K t03-c2-12-1 p
12. --- \VUgras{Mugės aikštėje} \textsuperscript{(39)} gausybė žmonių. Ateinama
pažiūrėti įvairių \VUgras{balaganų}: medinių \VUgras{arkliukų}
\textsuperscript{(40)}, \VUgras{cirko} \textsuperscript{(41)}, kuriame vaidinimas jau
prasidėjo; \VUgras{viešojo baliaus} \textsuperscript{(42)}, kur miesto jaunuoliai ir
jaunuolės mėgaujasi \VUgras{šokių} malonumu.

%%K t03-c2-13-1 p
13. --- Gale matome \VUgras{areną} \textsuperscript{(44)}, kurioje narsus ispanų
\VUgras{matadoras} kaunasi su įtūžusiu \VUgras{jaučiu} \textsuperscript{(45)}, paskui
\VUgras{kalnelius} \textsuperscript{(49)}, \VUgras{šaudyklą} \textsuperscript{(46)},
kur mankštinamasi vartoti \VUgras{šaunamuosius ginklus} ir šaudyti į
\VUgras{taikinį}.

%%K t03-c2-14-1 p
14. --- Netoliese štai \VUgras{mugės teatras} \textsuperscript{(47)}, kuriame
vaidinama \VUgras{komedija} ir \VUgras{drama}. Visai arti yra \VUgras{milžino}
\textsuperscript{(48)} ir \VUgras{nykštuko} balaganas. Pirmojo \VUgras{ūgis} yra du
metrai penkiolika centimetrų; antrasis vos tokio didumo kaip vaikas.

%%K t03-c2-15-1 p
15. --- Galiausiai štai didysis \VUgras{žvėrynas} \textsuperscript{(50)} su savo
\VUgras{laukiniais žvėrimis,} su savo \VUgras{tigru} \textsuperscript{(51)}, turinčiu
galingus \VUgras{nagus,} su savo \VUgras{liūtu} \textsuperscript{(52)}, turinčiu
tankius \VUgras{karčius,} su savo dviem \VUgras{žirafomis} \textsuperscript{(53)} ir
su savo \VUgras{drambliu} \textsuperscript{(54)}, kuris taip vikriai naudojasi savo
\VUgras{straubliu.}

%%K t03-c2-16-1 p
16. --- \VUgras{Dresuotojas} \textsuperscript{(55)}, kuris dresiruoja tuos žvėris,
stovi prie balagano durų.

\VUsaut{6.0mm}
\VUfilet{22mm}
